Une maladie contagieuse se répand chez les bovins.

Afin de détecter cette maladie, un laboratoire élabore un test de dépistage.

Les essais de ce test sur un groupe témoin de \np{1000} bovins sont résumés dans le tableau suivant :

\begin{center}
	\begin{tblr}{hlines, vlines, colspec={l*{3}{X[c,40mm]}}}
		& Test positif à la maladie & Test négatif à la maladie & Total \\
		Malade & 180 & 20 & 200 \\
		Non malade & 45 & 755 & 800 \\
		Total & 225 & 775 & \np{1000} \\
	\end{tblr}
\end{center}

On choisit au hasard un bovin ayant participé à l'essai et on note :
\begin{itemize}
	\item $M$ l'évènement \og le bovin est malade \fg{} ;
	\item $T$ l'évènement \og le test est positif \fg.
\end{itemize}

On exprimera tous les résultats sous forme décimale.

\begin{minipage}{0.7\linewidth}
	\begin{enumerate}
		\item Montrer que la probabilité de l'évènement $M$ est $0,2$.
		\item Décrire par une phrase l'évènement $M \cap T$, puis calculer la probabilité de cet évènement.
		\item Calculer la probabilité $P_M(T)$.
	\end{enumerate}
\end{minipage}
\hfill
\fbox{
	\begin{minipage}{0.2\textwidth}
		\textbf{Aide au calcul :}\\[0.5em]
		$225 = 5 \times 45$
	\end{minipage}
}
\begin{enumerate}[start=4]
	\item Un test est dit fiable si la probabilité de ne pas être atteint par la maladie sachant que le test est positif est inférieure ou égale à $0,15$.

	En se basant sur les résultats de cet essai, le test du laboratoire est-il fiable ? Justifier la réponse.
	\item À présent, on choisit deux bovins parmi ceux ayant participé à l'essai. Le nombre de bovins est suffisamment grand pour considérer cette expérience comme un tirage avec remise. On souhaite déterminer la probabilité de l'évènement $A$ : \og les deux bovins sont malades \fg.
	\begin{enumerate}
		\item Représenter la situation par un arbre pondéré.
		\item En déduire la probabilité de l'évènement $A$.
	\end{enumerate}
\end{enumerate}
